\documentclass{beamer}

\newcommand{\week}{GFC 5 of 5}

\title{Recovery and Aftermath}
\subtitle{Great Financial Crisis (GFC) of 2008}

\author{Econ 357}
\date{\week}

% Reference the shared preamble
\input{preamble}

% Set the footer -- change 
\setbeamertemplate{footline}{
  \leavevmode%
  \vspace{2ex}
  \hbox{%
    % Left box: Econ 457
    \begin{beamercolorbox}[wd=.4\paperwidth,ht=2.5ex,dp=1ex,left]{author in head/foot}%
      \hspace{1em}Econ 357
    \end{beamercolorbox}%
    % Middle box: Week
    \begin{beamercolorbox}[wd=.2\paperwidth,ht=2.5ex,dp=1ex,center]{date in head/foot}%
      \centering\week
    \end{beamercolorbox}%
    % Right box: Slide numbers
    \begin{beamercolorbox}[wd=.4\paperwidth,ht=2.5ex,dp=1ex,center]{date in head/foot}%
      \hfill\insertframenumber{} 
    \end{beamercolorbox}%
  }%
  \vskip0pt%
}

\begin{document}

\frame{\titlepage}

\begin{frame}
    \frametitle{Outline for Lectures on GFC of 2008}

    \begin{enumerate}
        \item Housing
        \item Leverage and Banks
        \item Bank Run
        \item Response
        \item \fbox{Recovery}
    \end{enumerate}
    \vspace{1em}

\end{frame}

\begin{frame}
    \frametitle{Outline for Today's Lecture}
    
\begin{enumerate}
    \item Highlights of Econ 357
    \item Recession - Causes
    \item Recession - Response
    \begin{itemize}
        \item Fiscal Policy
    \end{itemize}
    \item Reform
    \item Populist Backlash
    \begin{itemize}
        \item Tea Party and the election of 2010
    \end{itemize}
\end{enumerate}

\end{frame}

\begin{frame}
    \frametitle{1. Highlights of Econ 357}
    \section{Highlights of Econ 357}

    The following slides summarize what are, in my opinion, the coolest ideas we discussed in Econ 357.   To be clear, they are not meant to be an exhaustive review for the final.   Instead, they are simply what I find most interesting and what I hope you learned.

\end{frame}

\begin{frame}[t]
    \frametitle{1. Highlights of Econ 357}
    \framesubtitle{Financial Markets}

    Financial Markets
    \begin{itemize}
        \item Total returns = income plus capital gains
        \item Bond yields and prices are inversely related
        \item Stocks can be valued by dividends or P/E ratios, the later are high today.
    \end{itemize}

\end{frame}

\begin{frame}[t]
    \frametitle{1. Highlights of Econ 357}
    \framesubtitle{Financial Markets}

    \vspace{1em}

    \begin{columns}
    \begin{column}{0.55\textwidth}
    \centering
    \includegraphics[width=\textwidth]{Econ357/figures/02_mkts_2_bondprices.png}
    \end{column}
    \begin{column}{0.55\textwidth}
        \includegraphics[width=\textwidth]{Econ357/figures/02_mkts_7_pe_cape.png}
    \end{column}
    \end{columns}

    \blfootnote{See Markets Lectures \#2 and \#7}

\end{frame}

\begin{frame}[t]
    \frametitle{1. Highlights of Econ 357}
    \framesubtitle{Banks}

    \begin{itemize}
        \item Debt financing (i.e. leverage) magnifies positive returns, but also increases chances of failure
        \item Banks have high ROE because of leverage
        \item Higher capital makes banks safer (regulators like this) and reduces profits (investors don’t like this)
        \item Bank runs are a possible equilibrium.   Deposit insurance can prevent bank runs.
    \end{itemize}

\end{frame}

\begin{frame}[t]
    \frametitle{1. Highlights of Econ 357}
    \framesubtitle{Banks}

    \centering
    \includegraphics[width=0.9\textwidth]{Econ357/figures/03_banks_01_20pct.png}

    \blfootnote{See Banks, lecture \#1}

\end{frame}

\begin{frame}[t]
    \frametitle{1. Highlights of Econ 357}
    \framesubtitle{Banks}

  \centering
    \includegraphics[width=0.8\textwidth]{Econ357/figures/03_banks_5_bank_failures.png}

    \blfootnote{See Banks, lecture \#6}

\end{frame}



\begin{frame}[t]
    \frametitle{1. Highlights of Econ 357}
    \framesubtitle{Federal Reserve}

    \begin{itemize}
        \item Fed’s dual mandate is sometimes in conflict (like today)
        \item The Fed has lots of tools to adjust policy – interest rates (forward guidance) and balance sheet (loans and purchases)
    \end{itemize}

\end{frame}

\begin{frame}[t]
    \frametitle{1. Highlights of Econ 357}
    \framesubtitle{Federal Reserve}

      \centering
    \includegraphics[width=\textwidth]{Econ357/figures/05_gfc_5_fedfunds.png}

\end{frame}

\begin{frame}[t]
    \frametitle{1. Highlights of Econ 357}
    \framesubtitle{GFC of 2008}

    \begin{itemize}
        \item Mortgages use leverage, amplifying returns, both up and down
        \item Bank runs can happen in non-banks (i.e. shadow banks), and they did a lot in 2008.
        \item Events are rarely as obvious as they later seem to have been.
    \end{itemize}
    
\end{frame}

\begin{frame}[t]
    \frametitle{1. Highlights of Econ 357}
    \framesubtitle{GFC of 2008}
    
    \vspace{1em}

    \begin{columns}
    \begin{column}{0.55\textwidth}
    \centering
    \includegraphics[width=\textwidth]{Econ357/figures/05_gfc_3_case_shiller.png}
    \end{column}
    \begin{column}{0.55\textwidth}
        \includegraphics[width=\textwidth]{Econ357/figures/05_gfc_2_householdleverage.png}
    \end{column}
    \end{columns}
    
    \blfootnote{See GFC lectures \#1 and \#2}

\end{frame}

\begin{frame}[t]
    \frametitle{1. Highlights of Econ 357}
    \framesubtitle{GFC of 2008}

     \centering
     \includegraphics[width=\textwidth]{Econ357/figures/05_gfc_3_lehman.png}

\end{frame}

\begin{frame}[t]
    \frametitle{2. Recession - Causes}
    \section{Recession - Causes}


    While the panic phase subsided by early 2009, the aftermath lasted many years.   This period has been called the “Great Recession.”
    \vspace{1em}
    
    The Great Recession was characterized by the following:
    \begin{itemize}
        \item Credit contraction
        \item Debt overhang
        \item Sectoral overhang:
        \begin{itemize}
            \item Housing 
            \item Autos
        \end{itemize}
    \end{itemize}

\end{frame}

\begin{frame}[t]
    \frametitle{2. Recession - Causes}
    \framesubtitle{Credit Contraction}

    Banks pulled back sharply on lending, making it much harder to get credit to finance investments or consumption.
    \vspace{1em}

    \centering
    \includegraphics[width=0.8\textwidth]{Econ357/figures/05_gfc_5_causes1.png}
\end{frame}

\begin{frame}[t]
    \frametitle{2. Recession - Causes}
    \framesubtitle{Debt Overhang}

    During the 2000s household consumption was supported by increasing debt.   During the aftermath of the Great Recession, households borrowed less and consumed less (”deleveraging”)
   \vspace{1em}

    \centering
    \includegraphics[width=\textwidth]{Econ357/figures/05_gfc_5_causes2.png}

\end{frame}

\begin{frame}[t]
    \frametitle{2. Recession - Causes}
    \framesubtitle{Sectoral Overhang}

    Housing was especially hard hit for two reasons: (1) buying a new home requires credit, which was tight, and lower prices meant many homeowners had little to no equity. 

   \vspace{1em}

    \centering
    \includegraphics[width=\textwidth]{Econ357/figures/05_gfc_5_causes3.png}

\end{frame}

\begin{frame}[t]
    \frametitle{2. Recession - Causes}
    \framesubtitle{Sectoral Overhang}

   The auto sector was also especially hard hit.  Like buying a house, most consumers need credit to purchase a car.  GM received a “bridge loan” from the Federal Government, but eventually went through bankruptcy
   \vspace{1em}

    \centering
    \includegraphics[width=\textwidth]{Econ357/figures/05_gfc_5_causes4.png}

\end{frame}

\begin{frame}[t]
    \frametitle{2. Recession - Causes}
    \framesubtitle{Sectoral Overhang}

    \centering
    \includegraphics[width=\textwidth]{Econ357/figures/05_gfc_5_causes5.png}

\end{frame}

\begin{frame}[t]
    \frametitle{3. Recession - Response}
    \section{Recession - Response}

     Stopping the panic required:
    \begin{itemize}
        \item Recapitalize the financial sector, including by encouraging or forcing an increase in capital
        \item Supporting asset prices.
    \end{itemize}
    \vspace{1em}

    Supporting the recovery was the next phase.   This required more traditional counter-cyclical fiscal policy (i.e. when private households and businesses aren't spending, the government steps in and spends)



\end{frame}

\begin{frame}[t]
    \frametitle{3. Recession - Response}
    \framesubtitle{Recovery Act}

    \centering
    \includegraphics[width=\textwidth]{Econ357/figures/05_gfc_5_response1.png}

\end{frame}

\begin{frame}[t]
    \frametitle{3. Recession - Response}
    \framesubtitle{Recovery Act}
    
    American Reinvestment and Recovery Act (ARRA) of 2008
    
    Total = \$787 billion
    \begin{itemize}
        \item Tax Relief
        \item State and local fiscal relief
        \item Infrastructure
        \item Unemployment benefits / SNAP
    \end{itemize}
    \vspace{1em}
    
     The Congressional Budget Office estimated that ARRA had the following effects:  Raised GDP by 0.7 to 1.8\% in 2009, and 0.7 to 4.1\% in 2010
        Lowered the unemployment rate by 0.4 to 1.8\% by 2010.   \textit{All relative to a counterfactual of no government support.}
    
\end{frame}



\begin{frame}[t]
    \frametitle{3. Recession - Response}
    \framesubtitle{Recovery Act}

    \centering
    \includegraphics[width=\textwidth]{Econ357/figures/05_gfc_5_response2.png}

\end{frame}

\begin{frame}[t]
    \frametitle{4. Reform}
    \section{Reform}
    \framesubtitle{Dodd-Frank}

    The Dodd-Frank Wall Street Reform Act was passed in 2010.   The legislation included the following parts:
    \begin{itemize}
        \item Consumer Financial Protection Bureau (CFPB)
        \item Systemically Important Financial Institutions (SIFIs)
        \item Fed’s Emergency (13(3)) Loans
        \item Orderly Liquidation Authority (OLA)
        \item Stress Tests and Basel III
    \end{itemize}

\end{frame}

\begin{frame}[t]
    \frametitle{4. Reform}
    \framesubtitle{Consumer Financial Protection Bureau (CFPB}

    \textbf{Consumer Financial Protection Bureau (CFBP)}

    \begin{itemize}
        \item The SEC was created by FDR during the 1930s in order to rebuild trust in financial markets and to protect investors.
        \item In a similar way, the CFBP was created in 2010 in order to protect consumers of financial products.   Its authority has been disputed and recently it objectives have been dramatically scaled back.
    \end{itemize}

\end{frame}

\begin{frame}[t]
    \frametitle{4. Reform}
    \framesubtitle{Systemically Important Financial Institutions (SIFIs)}

    \textbf{Systemically Important Financial Institutions (SIFIs)}

    \begin{itemize}
        \item Prior to 2008 there were significant gaps in the regulatory system.   Bear Stearns was regulated by the SEC.   AIG didn’t really have a regulator for the financial holding company.
        \item SIFIs are now regulated by the Federal Reserve, which can impose capital requirements. 
        \begin{itemize} 
            \item AIG was determined to be a SIFI in 2013, but this was rescinded in 2017.
        \end{itemize}
    \end{itemize}

\end{frame}

\begin{frame}[t]
    \frametitle{4. Reform}
    \framesubtitle{Fed’s Emergency (13(3)) Loans}

    \textbf{Fed’s Emergency (13(3)) Loans}

    \begin{itemize}
        \item Prior to 2008 there were very few limitations on the Fed’s ability to make emergency loans.   The Fed used emergency loans for Bear Stearns and for AIG.
        \item The Fed is now restricted from making loans to single companies.   Emergency loans must be broad based and accessible to a wide range of financial market participants.  (In practice, this hasn't prevented the Fed from making 13(3) loans.  It made significant 13(3) loans in 2020 and again in 2023)
    \end{itemize}

\end{frame}

\begin{frame}[t]
    \frametitle{4. Reform}
    \framesubtitle{Orderly Liquidation Authority}

    \textbf{Orderly Liquidation Authority}

    \begin{itemize}
        \item Prior to 2008 it was very difficult to put any non-bank through bankruptcy.   (It took 10 years to resolve all the claims in the Lehman bankruptcy).
        \item The FDIC can now liquidate non-bank financial institutions.   This could help to make bankruptcy faster, and therefore a more accessible tool for regulators.
    \end{itemize}

\end{frame}

\begin{frame}[t]
    \frametitle{4. Reform}
    \framesubtitle{Stress Tests and Basel III}

    \textbf{Stress Tests and Basel III}
    \textit{Basel III was not actually part of Dodd-Frank, but is relevant here.}

    \begin{itemize}
        \item Prior to 2008 banks were able to use substantial leverage through a number of different tools (for example, off-balance sheet exposure to SIVs).
        \item Capital requirements are now substantially higher.   Banks undergo annual stress tests and must meet the proscribed capital levels.   Basel III is an international agreement on bank capital practices, and includes a capital surcharge for systemically important banks.
    \end{itemize}

\end{frame}

\begin{frame}[t]
    \frametitle{4. Reform}
    \framesubtitle{Stress Tests and Basel III}

    \centering
    \includegraphics[width=\textwidth]{Econ357/figures/05_gfc_5_reform.png}

\end{frame}

\begin{frame}[t]
    \frametitle{5. Populist Backlash}
    \framesubtitle{}

    In March 2009 AIG paid executives in its Financial Products division \$165 million in bonuses.   This caused a tremendous outrage:
    \vspace{1em}

    \myquote{If there is a single episode in this entire 18 months that has made me more angry, I can’t think of one other than AIG.}{
    Ben Bernanke \\
    Fed Chairman \\
    2009}

    
\end{frame}

\begin{frame}[t]
    \frametitle{5. Populist Backlash}
    \framesubtitle{}

    \myquote{Paying the creeps, I don’t care on what level they are on, is nuts…‘Let ’em go to court. Let ’em sue. Let ’em go to h***!}{
    LA Times reader \\
    (as quoted by Pew Research) \\
    2009}

\end{frame}

\begin{frame}[t]
    \frametitle{5. Populist Backlash}
    \framesubtitle{Tea Party Election}

    \textbf{Tea Party Election of 2010}
    \begin{itemize}
        \item In November 2010 Republicans gained 63 seats in the House of Representatives.
        \item A major contributor to the Republican success was anger about the Democrat’s response to the financial crisis.   Many of the newly elected representatives self-identified with the “tea party” wing of the Republican part.
        \item President Obama admitted that he “humbled” by the “shellacking”
    \end{itemize}

\end{frame}

\begin{frame}[t]
    \frametitle{5. Populist Backlash}
    \framesubtitle{Tea Party Election}

    \centering
    \includegraphics[width=0.9\textwidth]{Econ357/figures/05_gfc_5_teaparty.png}

\end{frame}


\end{document}